% RESEARCH_DRAFT_SHA256: aa732814039b095ef99c8dc71c78416ff24d086ca08aaede8a11beef36888718
% Layout and source-reference wrappers expanded for the web renderer.
\documentclass[11pt]{article}
\usepackage[margin=25mm]{geometry}
\usepackage{fontspec,unicode-math}
\setmainfont{Latin Modern Roman}
\setmathfont{Latin Modern Math}
\usepackage{amsmath,amssymb,mathtools,xurl,hyperref}
\hypersetup{colorlinks=true,urlcolor=blue,linkcolor=blue}
\setcounter{secnumdepth}{0}
\setlength{\parindent}{0pt}
\setlength{\parskip}{5pt}
\setlength{\emergencystretch}{3em}
\newcommand{\sref}[2]{\hyperlink{source:#1:#2}{[S#2]}}
\newcommand{\cataloguescope}[1]{\paragraph{Source record.}#1}
\newcommand{\floor}[1]{\left\lfloor#1\right\rfloor}
\begin{document}
\section*{Research-attempt provenance}
This standalone research attempt retains the complete original definition body below,
including its historical source and preparation notes. Those notes describe the
earlier catalogue preparation; the new attempt and its verification are identified
separately after that body. No original ID, catalogue status, or source statement is changed.

% UnsolvedMath ID 30005518; source catalogue code OWR-13750331-002
\section{Reconstruction of Chromosome Structure from Contact Maps}\label{30005518}
{\small UnsolvedMath ID 30005518 · Computer Science}
\subsection{Definitions and mathematical statement}

\paragraph{Shared notation.}Standard number systems are $\mathbb N,\mathbb Z,\mathbb Q,\mathbb R,\mathbb C$. Specialized notation and index ranges follow the cited source and the definitions below.


A contact map records detected proximity between pairs of sites. A distance matrix records proposed pairwise distances. A Gram matrix contains inner products and can reconstruct coordinates up to rigid transformations when its distance data are consistent.
\paragraph{Record scope and qualifications.}
The catalogue title is retained as metadata. Named certificates, contracts, bounds or frameworks in that title are not asserted as results of the original source extract.
The supplied registry flags this extraction as unrecoverable. It is displayed as a literal source fragment; the definitions provide context rather than a replacement formal statement.

The specific contact model, reconstruction criterion, and separate research question are not recovered from the mixed explanatory text.

Source-record limitation: Source flags the statement as unrecoverable: The extracted phrase is a figure caption answered by the pictured contact map; the surrounding contribution describes established three-dimensional reconstruction methods rather than an open problem. No separate clean mathematical statement is supplied; this source ID preserves its provenance.
The source text below is rendered literally where necessary; damaged formulas and composite extracts are not silently repaired into a new mathematical claim.
The supplied catalogue marks this record as unrecoverable. Its extract and terminology are documented, but a complete original problem has not been reconstructed.
\paragraph{Statement recorded in the supplied source.}
What does a chromosome contact map look like? represent and analyze the structure of proteins. They are closely related to re- currence networks obtained from recurrence plots of dynamical systems, and the same reconstruction methods can be used. The main steps of the reconstruction are (i) to derive a complete distance matrix using shortest-path distance on the weighted contact network, (ii) to use a result of distance geometry based on the three dominant eigenvectors of the metric matrix (yielding the Gram matrix of the structure) (iii) or approximating the 3D structure using multidimensional scaling techniques. Analogies between the algebraic analysis (principal eigenvector of the correla- tion matrix) and motif analysis (topological domains seen as squares along the diagonal) of the corresponding contact matrices can also be drawn. \href{https://doi.org/10.4171/owr/2023/20}{[S2]}
\paragraph{Source status.}
The supplied catalogue labels this source record partially solved; that metadata was checked against the supplied snapshot on 2026-10-11. The label is preserved for provenance and is not an independent verification of present open status or mathematical correctness. The catalogue literature assessment is dated 2026-08-22.
\subsection{Short English statement}
The fragment discusses chromosome or protein geometry reconstructed from contact information.
\subsection{Sources}
\begin{itemize}
\item[S1] Ulam AI and the UnsolvedMath Contributors, UnsolvedMath: A Curated Collection of Open Mathematics Problems, record 30005518 (OWR-13750331-002). Catalogue curation and metadata: CC BY 4.0. The source extract is retained; definitions, English summaries and qualifications are editorial additions.
\url{https://huggingface.co/datasets/ulamai/UnsolvedMath}.
\item[S2] Mathematical Foundations of Biological Organisation (2023), Oberwolfach Reports, DOI 10.4171/owr/2023/20, pp. 1085--1104
\url{https://doi.org/10.4171/owr/2023/20}.
\item[S3] Hi-C primary study
\url{https://doi.org/10.1016/j.cell.2014.11.021}.
\item[S4] multicell-type contact maps
\url{https://doi.org/10.1016/j.celrep.2016.10.061}.
\item Source-specific Oberwolfach Reports licences, including CC BY-SA 4.0 where applicable, are retained; see THIRD\_PARTY\_NOTICES.md.
\end{itemize}
\paragraph{Preparation and review.}
The original source and UnsolvedMath contributors retain attribution for the source extract. Codex (GPT-6; exact serving revision not exposed) assisted with definitions, English summaries, scope notes and formatting. No new human mathematical authorship or independently verified solution is asserted. Permanent source ID 30005518 and code OWR-13750331-002 checked against the complete supplied registry and current repository. Distinct upstream records, including repeated questions and broad or damaged extracts, are retained for complete coverage of the source CS category. A repeated record is not a new independent conjecture. Complete mathematical-equivalence and current-literature review has not been performed. Preparation and final review follow contribution prompts C and D, with source gaps disclosed. Structural and rendering checks do not establish a new theorem or open status.
\label{end:30005518}

\section{Literature and context check}

I read Lesne's contribution on pages 1093--1094 of the primary report.
The opening question in the extract comes from a figure caption; the surrounding
text describes a reconstruction pipeline and approximation by multidimensional
scaling, not a precise uniqueness theorem. The calculation below tests an exact
inference one might mistakenly make from that pipeline. It does not interpret
a contact frequency directly as a Euclidean distance or assert a biological model.

\section{Research attempt: An exact Euclidean-distance test and a shortest-path counterexample}

\paragraph{When a complete distance array admits an exact embedding.}
Let $D=(d_{ij})$ be a symmetric nonnegative $N\times N$ array with zero
diagonal, $N\geq1$. Put $S_{ij}=d_{ij}^2$,
$H=I-\mathbf1\mathbf1^T/N$, and $B=-\frac12HSH$.
Then $D$ is the distance array of points in a Euclidean space if and only if
$B\succeq0$. The least embedding dimension is $\operatorname{rank}B$.
For the forward implication, center the points so their coordinate matrix $Z$
has $Z^T\mathbf1=0$. Their squared distances are
$S_{ij}=\|z_i\|^2+\|z_j\|^2-2z_i^Tz_j$, whence $B=ZZ^T\succeq0$.
For the converse, factor the PSD matrix $B=ZZ^T$ with $\operatorname{rank}B$
columns by its nonnegative eigenvalues. Since $B\mathbf1=0$, one has
$\|Z^T\mathbf1\|^2=\mathbf1^TB\mathbf1=0$, so the points are centered.
Let $S'$ be their squared-distance matrix. Then $HS'H=-2B=HSH$.
Writing $E=S-S'$, the equality $HEH=0$ means
$E_{ij}=r_i+r_j-r$, where $r_i$ is its $i$th row mean and $r$ is the overall
mean. Its zero diagonal gives $2r_i=r$ for every $i$. Averaging gives
$2r=r$, hence every $r_i=0$ and $E=0$. Thus the constructed distances are
exactly $D$. Any embedding has centered Gram matrix $B$ and needs at least
its rank many coordinates, proving the dimension claim.

\paragraph{A shortest-path metric that fails the test.}
Take the unweighted four-cycle. Its shortest-path distances are one between
adjacent vertices and two between opposite vertices, so
\[
 S=\begin{pmatrix}0&1&4&1\\1&0&1&4\\4&1&0&1\\1&4&1&0\end{pmatrix}.
\]
The vector $v=(1,-1,1,-1)^T$ has sum zero and $Sv=2v$.
Therefore $Bv=-v$ and $v^TBv=-4<0$.
Although shortest-path distances satisfy symmetry, nonnegativity and the triangle
inequality, they are not Euclidean in any dimension in this example.
The remaining centered eigenvectors $(1,0,-1,0)$ and $(0,1,0,-1)$ have
$S$-eigenvalue $-4$, hence $B$-eigenvalue $2$; the constant vector has
$B$-eigenvalue zero. The spectrum is exactly $2,2,0,-1$.
Keeping the positive eigenspaces gives points
$(1,0),(0,1),(-1,0),(0,-1)$. Their opposite distances are two but adjacent
distances are $\sqrt2$, not one. Spectral truncation has produced an
approximation with explicitly changed measurements, not exact recovery.

\paragraph{Even a PSD matrix can need more than three coordinates.}
If $N=5$ and all off-diagonal distances are one, then $S=J-I$ and
$B=\frac12H$ has rank four, as $H$ is identity on the four-dimensional
zero-sum subspace. It is a legitimate Euclidean distance array but cannot be
represented exactly in three dimensions. Thus neither taking three dominant
eigenvectors nor taking all nonnegative eigenvalues universally justifies an
exact 3D reconstruction claim. A source-specific noise objective and dimension
assumption are needed to evaluate the approximation.

\section{Verification and adversarial audit}

The squared-distance array, rather than the distance array itself, is centered.
The four-cycle failure has an exact negative quadratic-form witness and all
four eigenvalues are checked symbolically. The regular-simplex example separates
PSD feasibility from the rank-three requirement. The $N=1$ case has zero
Gram matrix and embeds as a single point.

\paragraph{Audit corrections.} Identified the figure-caption provenance and distinguished graph metrics,
Euclidean metrics and rank-three embeddings. Computed the distortion introduced
by the positive spectral truncation instead of calling it exact recovery.
\paragraph{External computation.} None; symbolic arguments and boundary checks only.
\section{Scope and remaining gaps}

This is a complete diagnostic for supplied complete distances and two concrete
obstructions to exact use of the described pipeline. The original record lacks
a definite inverse problem, contact-to-distance law or uniqueness criterion.
The calculation does not characterize chromosome maps or supply a validated
reconstruction method for their biological measurements.

\paragraph{Integration review.} Codex (GPT-6; exact serving revision not exposed) reviewed the written arguments and expanded layout and source-reference wrappers. Repository metadata and the production web renderer were checked using Python 3.13 and JavaScript with MathJax 3.2.2. This model-assisted review is not independent human proof verification; no PDF compilation or formal proof checking is asserted.

\section{Completion Estimate}
\noindent\textbf{COMPLETION: 0\%}
The calculations below the preserved source statement investigate a concrete route,
but establish no substantiated new progress toward the entire original target.
Known special cases and failed sufficient conditions are not credited as a solution.
\section{References}
The original source references and catalogue attribution are preserved in the
Sources section above. The following URLs were actually consulted for this attempt;
retention of an additional original reference does not claim fresh consultation of it.
\begin{itemize}
\item \url{https://ems.press/content/serial-article-files/47015} (consulted 2026-10-11). 
\end{itemize}
\section{Model and contributor disclosure}
Codex GPT-6 completed written-proof review of this attempt at its declared scope; this is model-assisted review, not independent human proof checking.
OpenAI gpt-6-astra, requested reasoning configuration ultra, generated this research
attempt on 2026-10-11. The exact serving revision is not exposed. Codex GPT-6
performed orchestration and repository-format review; its exact serving revision
is also not exposed. The mathematical route, boundary checks and final adversarial
self-audit followed contribution prompts A then B in
\texttt{docs/contribution\_prompts.md}, after reading \texttt{CONTRIBUTING.md}.
Python 3.14.2 assembled the standalone documents; any mathematical computation is
specified in Verification. No theorem prover was used. No human mathematical
authorship, substantial human proof correction, or independent human proof checking
is asserted. Model-assisted review, source retrieval, and structural checks do not
establish correctness or novelty. Full private reasoning transcripts are unavailable
as reproducibility artifacts; the complete arguments, inputs and outcomes actually
used are stated here or in the companion batch audit.
\section{Final status}
Unresolved as a new research contribution. This is a concrete route investigation,
not a claim of new complete resolution or a current-status certification. A known
answer reproduced for context does not become a new contribution. The catalogue's
source-status label is preserved; this attempt status does not assert that a
source-labelled solved question is open.
\end{document}
